% ============================================================================
% CAPÍTULO 2: INTRODUÇÃO À MODELAGEM MATEMÁTICA
% Bioinformática para Biologia de Sistemas
% ============================================================================

% ============================================================================
% TEMPLATE BASE PARA APRESENTAÇÕES EM BEAMER
% Bioinformática para Biologia de Sistemas
% ============================================================================

\documentclass[12pt, aspectratio=169]{beamer}

% ============================================================================
% PACOTES ESSENCIAIS
% ============================================================================
\usepackage[utf8]{inputenc}
\usepackage[T1]{fontenc}
\usepackage[brazilian]{babel}
\usepackage{amsmath, amssymb, amsthm}
\usepackage{graphicx}
\usepackage{xcolor}
\usepackage{tikz}
\usetikzlibrary{arrows.meta, shapes, positioning, calc, decorations.pathreplacing, decorations.shapes, decorations.pathmorphing}
\usepackage{listings}
\usepackage{booktabs}
\usepackage{multirow}
\usepackage{hyperref}
\usepackage{chemfig}
\usepackage{chemarr}

% ============================================================================
% CONFIGURAÇÃO DO TEMA
% ============================================================================
\usetheme{Madrid}
\usecolortheme{default}

% Cores personalizadas
\definecolor{azulescuro}{RGB}{0, 51, 102}
\definecolor{azulclaro}{RGB}{51, 153, 255}
\definecolor{verdeacido}{RGB}{102, 204, 0}
\definecolor{laranjacel}{RGB}{255, 153, 51}
\definecolor{roxodna}{RGB}{153, 51, 255}

\setbeamercolor{structure}{fg=azulescuro}
\setbeamercolor{palette primary}{bg=azulescuro,fg=white}
\setbeamercolor{palette secondary}{bg=azulclaro,fg=white}
\setbeamercolor{palette tertiary}{bg=azulescuro,fg=white}
\setbeamercolor{palette quaternary}{bg=azulclaro,fg=white}
\setbeamercolor{block title}{bg=azulescuro,fg=white}
\setbeamercolor{block body}{bg=azulclaro!10,fg=black}
\setbeamercolor{block title alerted}{bg=red!80,fg=white}
\setbeamercolor{block body alerted}{bg=red!10,fg=black}
\setbeamercolor{block title example}{bg=verdeacido!80,fg=white}
\setbeamercolor{block body example}{bg=verdeacido!10,fg=black}

% ============================================================================
% CONFIGURAÇÃO DE CÓDIGO
% ============================================================================
\lstset{
    language=Python,
    basicstyle=\ttfamily\small,
    keywordstyle=\color{azulescuro}\bfseries,
    commentstyle=\color{verdeacido},
    stringstyle=\color{laranjacel},
    numbers=left,
    numberstyle=\tiny\color{gray},
    stepnumber=1,
    numbersep=5pt,
    backgroundcolor=\color{white},
    showspaces=false,
    showstringspaces=false,
    showtabs=false,
    frame=single,
    rulecolor=\color{black},
    tabsize=2,
    captionpos=b,
    breaklines=true,
    breakatwhitespace=false,
    escapeinside={\%*}{*)},
    xleftmargin=10pt,
    xrightmargin=10pt
}

% Definir estilo para R
\lstdefinestyle{Rstyle}{
    language=R,
    basicstyle=\ttfamily\small,
    keywordstyle=\color{azulescuro}\bfseries,
    commentstyle=\color{verdeacido},
    stringstyle=\color{laranjacel}
}

% Definir estilo para MATLAB
\lstdefinestyle{MATLABstyle}{
    language=Matlab,
    basicstyle=\ttfamily\small,
    keywordstyle=\color{azulescuro}\bfseries,
    commentstyle=\color{verdeacido},
    stringstyle=\color{laranjacel}
}

% ============================================================================
% COMANDOS PERSONALIZADOS
% ============================================================================

% Comando para equações destacadas
\newcommand{\destaque}[1]{%
    \begin{alertblock}{}
        \begin{center}
            #1
        \end{center}
    \end{alertblock}
}

% Comando para definições
\newcommand{\definicao}[2]{%
    \begin{block}{Definição: #1}
        #2
    \end{block}
}

% Comando para exemplos
\newcommand{\exemplo}[2]{%
    \begin{exampleblock}{Exemplo: #1}
        #2
    \end{exampleblock}
}

% Comando para observações importantes
\newcommand{\importante}[1]{%
    \begin{alertblock}{Importante!}
        #1
    \end{alertblock}
}

% Comandos matemáticos úteis
\newcommand{\R}{\mathbb{R}}
\newcommand{\N}{\mathbb{N}}
\newcommand{\Z}{\mathbb{Z}}
\newcommand{\dif}{\mathrm{d}}
\newcommand{\parcial}[2]{\frac{\partial #1}{\partial #2}}
\newcommand{\derivada}[2]{\frac{\dif #1}{\dif #2}}
\newcommand{\vect}[1]{\boldsymbol{#1}}

% Comandos biológicos
\newcommand{\gene}[1]{\textit{#1}}
\newcommand{\proteina}[1]{\textsc{#1}}
\newcommand{\especie}[1]{\textit{#1}}

% ============================================================================
% CONFIGURAÇÃO DE RODAPÉ
% ============================================================================
\setbeamertemplate{footline}{
  \leavevmode%
  \hbox{%
  \begin{beamercolorbox}[wd=.33\paperwidth,ht=2.25ex,dp=1ex,center]{author in head/foot}%
    \usebeamerfont{author in head/foot}\insertshortauthor
  \end{beamercolorbox}%
  \begin{beamercolorbox}[wd=.34\paperwidth,ht=2.25ex,dp=1ex,center]{title in head/foot}%
    \usebeamerfont{title in head/foot}\insertshorttitle
  \end{beamercolorbox}%
  \begin{beamercolorbox}[wd=.33\paperwidth,ht=2.25ex,dp=1ex,right]{date in head/foot}%
    \usebeamerfont{date in head/foot}\insertshortdate{}\hspace*{2em}
    \insertframenumber{} / \inserttotalframenumber\hspace*{2ex}
  \end{beamercolorbox}}%
  \vskip0pt%
}

% ============================================================================
% CONFIGURAÇÃO DE NAVEGAÇÃO
% ============================================================================
\setbeamertemplate{navigation symbols}{}

% ============================================================================
% CONFIGURAÇÃO DE ITEMIZE/ENUMERATE
% ============================================================================
\setbeamertemplate{itemize items}[circle]
\setbeamertemplate{enumerate items}[default]

% ============================================================================
% FIM DO TEMPLATE
% ============================================================================


% ============================================================================
% INFORMAÇÕES DO DOCUMENTO
% ============================================================================
\title[Cap. 2: Modelagem Matemática]{Capítulo 2: Introdução à Modelagem Matemática}
\subtitle{Bioinformática para Biologia de Sistemas}
\author{Ney Lemke}
\institute{DBF-Unesp}
\date{\today}

\usetikzlibrary{decorations.pathmorphing, positioning}
\begin{document}

% ============================================================================
% SLIDE DE TÍTULO
% ============================================================================
\begin{frame}
\titlepage
\end{frame}

% ============================================================================
% SUMÁRIO
% ============================================================================
\begin{frame}{Sumário}
\tableofcontents
\end{frame}

% ============================================================================
% SEÇÃO 1: INTRODUÇÃO
% ============================================================================
\section{Introdução}

\begin{frame}[shrink=5]{2.1 Introdução à Modelagem Matemática}
\begin{block}{Objetivos de Aprendizagem}
\begin{itemize}
    \item Compreender o que são modelos matemáticos
    \item Aprender a construir modelos para sistemas biológicos
    \item Dominar técnicas de análise de modelos
    \item Aplicar modelagem a problemas reais
\end{itemize}
\end{block}

\vspace{0.3cm}

\definicao{Modelo Matemático}{
Representação abstrata de um sistema real usando linguagem matemática, que permite fazer previsões quantitativas sobre o comportamento do sistema.
}
\end{frame}

\begin{frame}{Por que Modelar Sistemas Biológicos?}
\begin{columns}
\column{0.48\textwidth}
\textbf{Vantagens da Modelagem:}
\begin{itemize}
    \item Previsão de comportamento
    \item Teste de hipóteses \textit{in silico}
    \item Otimização de experimentos
    \item Identificação de componentes chave
    \item Design racional de sistemas
\end{itemize}

\column{0.48\textwidth}
\textbf{Aplicações:}
\begin{itemize}
    \item Dinâmica de populações
    \item Vias metabólicas
    \item Redes regulatórias
    \item Farmacocinética
    \item Biologia sintética
    \item Medicina personalizada
\end{itemize}
\end{columns}

\vspace{0.3cm}

\importante{
"All models are wrong, but some are useful" --- George Box
}
\end{frame}

\begin{frame}{O Ciclo da Modelagem}
\begin{center}
\begin{tikzpicture}[
    auto,
    excalidraw/.style={
        decorate,
        decoration={random steps, segment length=3pt, amplitude=0.5pt},
        thick,
        font=\fontfamily{pzc}\selectfont\large
    },
    box/.style={
        rectangle, 
        draw=azulescuro, 
        excalidraw,
        fill=azulclaro!5, 
        text width=2.4cm, 
        text centered, 
        minimum height=1.2cm
    },
    arrow/.style={
        ->, 
        thick, 
        azulescuro, 
        excalidraw
    }
]

    \node[box] (bio) {Sistema Biológico};
    \node[box, right=2cm of bio] (modelo) {Modelo Matemático};
    \node[box, below=1.5cm of modelo] (analise) {Analise Computacional};
    \node[box, left=2cm of analise] (pred) {Previsões\\Hipóteses};

    \draw[arrow] (bio) -- node[above, font=\footnotesize] {Abstração} (modelo);
    \draw[arrow] (modelo) -- node[right, font=\footnotesize] {Simulação} (analise);
    \draw[arrow] (analise) -- node[below, font=\footnotesize] {Interpretação} (pred);
    \draw[arrow] (pred) -- node[left, font=\footnotesize] {Validação} (bio);

\end{tikzpicture}
\end{center}

\begin{exampleblock}{Ciclo Iterativo}
O processo é iterativo: modelos são refinados com base em novos dados experimentais
\end{exampleblock}
\end{frame}

% ============================================================================
% SEÇÃO 2: TIPOS DE MODELOS
% ============================================================================
\section{Tipos de Modelos}

\begin{frame}{Classificação de Modelos}
\begin{block}{Por Natureza Matemática}
\begin{itemize}
    \item \textbf{Determinísticos:} resultado único para cada conjunto de condições
    \item \textbf{Estocásticos:} incorporam aleatoriedade e probabilidade
\end{itemize}
\end{block}

\begin{block}{Por Dependência Temporal}
\begin{itemize}
    \item \textbf{Estáticos:} não mudam com o tempo (redes, equilíbrio)
    \item \textbf{Dinâmicos:} evolução temporal (EDOs, EDPs)
\end{itemize}
\end{block}

\begin{block}{Por Escala}
\begin{itemize}
    \item \textbf{Microscópicos:} nível molecular/celular
    \item \textbf{Mesoscópicos:} populações celulares
    \item \textbf{Macroscópicos:} tecidos, órgãos, organismos
\end{itemize}
\end{block}
\end{frame}

\begin{frame}[shrink=10]{Modelos Contínuos vs Discretos}
\begin{columns}
\column{0.48\textwidth}
\begin{block}{Modelos Contínuos}
\begin{itemize}
    \item Variáveis contínuas
    \item Equações diferenciais
    \item Concentrações
    \item Grande número de moléculas
\end{itemize}

\textbf{Exemplo:}
\begin{equation*}
\derivada{[X]}{t} = k_1 - k_2[X]
\end{equation*}
\end{block}

\column{0.48\textwidth}
\begin{block}{Modelos Discretos}
\begin{itemize}
    \item Contagem de moléculas
    \item Algoritmos estocásticos
    \item Eventos discretos
    \item Pequeno número de moléculas
\end{itemize}

\textbf{Exemplo:}
\begin{equation*}
X(t+1) = X(t) + \text{prod} - \text{deg}
\end{equation*}
\end{block}
\end{columns}

% \vspace{0.3cm}

\begin{exampleblock}{Quando usar cada tipo?}
\begin{itemize}
    \item Contínuo: \textgreater 100 moléculas
    \item Discreto: \textless 100 moléculas
\end{itemize}
\end{exampleblock}
\end{frame}

\begin{frame}[shrink=20]{Modelos Empíricos vs Mecanísticos}
\small
\begin{columns}
\column{0.48\textwidth}
\begin{block}{Empíricos}
\textbf{Baseados em dados}

\begin{itemize}
    \item Ajuste de curvas
    \item Regressão estatística
    \item Machine learning
    \item Pouca interpretação biológica
\end{itemize}

\vspace{0.2cm}

\textbf{Vantagens:}
\begin{itemize}
    \item Simples e rápidos
    \item Bom ajuste aos dados
\end{itemize}

\textbf{Desvantagens:}
\begin{itemize}
    \item Pouco poder preditivo
    \item Extrapolação arriscada
\end{itemize}
\end{block}

\column{0.48\textwidth}
\begin{block}{Mecanísticos}
\textbf{Baseados em mecanismos}

\begin{itemize}
    \item Representam processos reais
    \item Leis de conservação
    \item Princípios físico-químicos
    \item Interpretação biológica
\end{itemize}

\vspace{0.2cm}

\textbf{Vantagens:}
\begin{itemize}
    \item Poder preditivo
    \item Extrapolação confiável
\end{itemize}

\textbf{Desvantagens:}
\begin{itemize}
    \item Complexos
    \item Muitos parâmetros
\end{itemize}
\end{block}
\end{columns}
\end{frame}

% ============================================================================
% SEÇÃO 3: CONSTRUÇÃO DE MODELOS
% ============================================================================
\section{Construção de Modelos}

\begin{frame}{Etapas da Construção de Modelos}
\begin{enumerate}
    \item \textbf{Definição do problema}
    \begin{itemize}
        \item Objetivo claro
        \item Escopo bem definido
        \item Questões a responder
    \end{itemize}

    \item \textbf{Identificação de variáveis}
    \begin{itemize}
        \item Variáveis de estado
        \item Parâmetros
        \item Entradas e saídas
    \end{itemize}

    \item \textbf{Formulação de equações}
    \begin{itemize}
        \item Leis de conservação
        \item Cinéticas de reação
        \item Condições iniciais e de contorno
    \end{itemize}

    \item \textbf{Analise e simulação}
\end{enumerate}
\end{frame}

\begin{frame}{Exemplo: Modelo de Decaimento Radioativo}
\begin{block}{Sistema}
Um isótopo radioativo decai espontaneamente ao longo do tempo
\end{block}

\textbf{1. Variável de estado:} $N(t)$ = número de núcleos no tempo $t$

\textbf{2. Hipótese:} Taxa de decaimento proporcional ao número de núcleos

\textbf{3. Equação:}
\destaque{
\begin{equation*}
\derivada{N}{t} = -\lambda N
\end{equation*}
}

\textbf{4. Solução analítica:}
\begin{equation*}
N(t) = N_0 e^{-\lambda t}
\end{equation*}

onde $N_0 = N(0)$ é a condição inicial e $\lambda$ é a constante de decaimento.
\end{frame}

\begin{frame}[fragile, allowframebreaks]{Implementação: Decaimento Radioativo}
\begin{lstlisting}[language=Python, caption=Solução numérica do decaimento, basicstyle=\ttfamily\footnotesize]
import numpy as np
import matplotlib.pyplot as plt
from scipy.integrate import odeint

# Definir a EDO
def decay(N, t, lam):
    """Taxa de variacao de N"""
    return -lam * N

# Parametros
lambda_val = 0.5  # constante de decaimento
N0 = 100          # quantidade inicial
t = np.linspace(0, 10, 100)
\end{lstlisting}
\end{frame}

\begin{frame}[fragile, shrink=5]{Implementação: Decaimento Radioativo (2/2)}
\begin{lstlisting}[language=Python, caption=Solução e Visualização]
# Resolver numericamente
N_num = odeint(decay, N0, t, args=(lambda_val,))

# Solucao analitica
N_ana = N0 * np.exp(-lambda_val * t)

# Plotar
plt.plot(t, N_num, 'b-', label='Numerica')
plt.plot(t, N_ana, 'r--', label='Analitica')
plt.xlabel('Tempo')
plt.ylabel('N(t)')
plt.legend()
\end{lstlisting}
\end{frame}

% ============================================================================
% SEÇÃO 4: EQUAÇÕES DIFERENCIAIS
% ============================================================================
\section{Equações Diferenciais em Biologia}

\begin{frame}[shrink=5]{Equações Diferenciais Ordinárias (EDOs)}
\definicao{EDO}{
Equação que relaciona uma função com suas derivadas em relação a uma variável independente (geralmente o tempo).
}

\vspace{0.3cm}

\begin{block}{Forma Geral}
\begin{equation*}
\derivada{\vect{x}}{t} = \vect{f}(\vect{x}, t, \vect{p})
\end{equation*}

onde:
\begin{itemize}
    \item $\vect{x}$ = vetor de variáveis de estado
    \item $t$ = tempo
    \item $\vect{p}$ = vetor de parâmetros
    \item $\vect{f}$ = função que define a dinâmica
\end{itemize}
\end{block}
\end{frame}

\begin{frame}[shrink]{Classificação de EDOs}
\begin{block}{Por Ordem}
\begin{itemize}
    \item \textbf{Primeira ordem:} $\derivada{x}{t} = f(x,t)$
    \item \textbf{Segunda ordem:} $\frac{d^2x}{dt^2} = f(x, \derivada{x}{t}, t)$
    \item \textbf{Ordem n:} derivadas até $n$-ésima ordem
\end{itemize}
\end{block}

\begin{block}{Por Linearidade}
\begin{itemize}
    \item \textbf{Linear:} $\derivada{x}{t} = ax + b$
    \item \textbf{Não-linear:} $\derivada{x}{t} = ax^2 + bx$
\end{itemize}
\end{block}

\begin{block}{Por Autonomia}
\begin{itemize}
    \item \textbf{Autônoma:} $\derivada{x}{t} = f(x)$ (não depende explicitamente de $t$)
    \item \textbf{Não-autônoma:} $\derivada{x}{t} = f(x,t)$
\end{itemize}
\end{block}
\end{frame}

\begin{frame}[shrink=10]{Exemplo: Crescimento Exponencial}
\begin{block}{Modelo Simples de Crescimento}
Uma populacao cresce proporcionalmente ao seu tamanho:
\begin{equation*}
\derivada{N}{t} = rN
\end{equation*}

onde $r$ é a taxa de crescimento per capita.
\end{block}

\textbf{Solução analítica:}
\begin{equation*}
N(t) = N_0 e^{rt}
\end{equation*}

\begin{columns}
\column{0.48\textwidth}
\textbf{Se $r > 0$:} crescimento\\
\textbf{Se $r < 0$:} decaimento\\
\textbf{Se $r = 0$:} estacionário

\column{0.48\textwidth}
\begin{exampleblock}{Tempo de Duplicação}
$t_{2} = \frac{\ln 2}{r}$
\end{exampleblock}
\end{columns}
\end{frame}

\begin{frame}[shrink=15]{Exemplo: Crescimento Logístico}
\begin{block}{Modelo de Verhulst}
Populacao com capacidade de suporte limitada $K$:
\begin{equation*}
\derivada{N}{t} = rN\left(1 - \frac{N}{K}\right)
\end{equation*}
\end{block}

\textbf{Solução analítica:}
\begin{equation*}
N(t) = \frac{K N_0 e^{rt}}{K + N_0(e^{rt} - 1)}
\end{equation*}

\small
\begin{columns}
\column{0.48\textwidth}
\textbf{Comportamento:}
\begin{itemize}
    \item $N \to K$ quando $t \to \infty$
    \item Forma sigmoidal
    \item Ponto de inflexão em $N = K/2$
\end{itemize}

\column{0.48\textwidth}
\begin{center}
\begin{tikzpicture}[scale=0.6]
    \draw[->] (0,0) -- (5,0) node[right] {$t$};
    \draw[->] (0,0) -- (0,4) node[above] {$N$};
    \draw[thick, azulescuro, domain=0:4.5, smooth, samples=100]
        plot (\x, {3/(1+2*exp(-1.5*\x))});
    \draw[dashed, verdeacido] (0,3) -- (5,3) node[right, font=\tiny] {$K$};
    \node[below] at (2.5,0) {Tempo};
\end{tikzpicture}
\end{center}
\end{columns}
\end{frame}

\begin{frame}[fragile, allowframebreaks]{Implementação: Crescimento Logístico}
\begin{lstlisting}[language=Python, caption=Modelo logístico, basicstyle=\ttfamily\footnotesize]
import numpy as np
import matplotlib.pyplot as plt
from scipy.integrate import odeint

def logistic(N, t, r, K):
    """Modelo de crescimento logistico"""
    return r * N * (1 - N/K)

# Parametros
r = 0.5    # taxa de crescimento
K = 100    # capacidade de suporte
N0 = 10    # populacao inicial

# Tempo
t = np.linspace(0, 20, 200)

# Resolver
N = odeint(logistic, N0, t, args=(r, K))
\end{lstlisting}
\end{frame}

\begin{frame}[fragile, shrink=5]{Implementação: Crescimento Logístico (2/2)}
\begin{lstlisting}[language=Python, caption=Visualização]
# Plotar
plt.figure(figsize=(10, 6))
plt.plot(t, N, 'b-', linewidth=2, label='N(t)')
plt.axhline(y=K, color='r', linestyle='--',
            label=f'K = {K}')
plt.xlabel('Tempo')
plt.ylabel('Populacao N(t)')
plt.legend()
plt.grid(True)
\end{lstlisting}
\end{frame}

% ============================================================================
% SEÇÃO 5: SISTEMAS DE EDOs
% ============================================================================
\section{Sistemas de EDOs}

\begin{frame}[shrink=15]{Sistemas Acoplados}
\begin{block}{Sistema Geral}
Quando múltiplas variáveis interagem:
\begin{align*}
\derivada{x_1}{t} &= f_1(x_1, x_2, \ldots, x_n, t)\\
\derivada{x_2}{t} &= f_2(x_1, x_2, \ldots, x_n, t)\\
&\vdots\\
\derivada{x_n}{t} &= f_n(x_1, x_2, \ldots, x_n, t)
\end{align*}
\end{block}

\textbf{Notação vetorial:}
\destaque{
\begin{equation*}
\derivada{\vect{x}}{t} = \vect{f}(\vect{x}, t)
\end{equation*}
}
\end{frame}

\begin{frame}[shrink=15]{Exemplo: Modelo Predador-Presa (Lotka-Volterra)}
\begin{block}{Sistema}
\small
\begin{align*}
\derivada{V}{t} &= \alpha V - \beta V P \quad \text{(Presas)}\\
\derivada{P}{t} &= \delta \beta V P - \gamma P \quad \text{(Predadores)}
\end{align*}

onde:
\begin{itemize}
    \item $V$ = populacao de presas
    \item $P$ = populacao de predadores
    \item $\alpha$ = taxa de crescimento das presas
    \item $\beta$ = taxa de predação
    \item $\gamma$ = taxa de morte dos predadores
    \item $\delta$ = eficiência de conversão
\end{itemize}
\end{block}
\end{frame}

\begin{frame}[shrink=10]{Interpretação do Modelo Predador-Presa}
\begin{columns}
\column{0.48\textwidth}
\textbf{Equação das Presas:}
\begin{equation*}
\derivada{V}{t} = \underbrace{\alpha V}_{\text{crescimento}} - \underbrace{\beta V P}_{\text{predação}}
\end{equation*}

\begin{itemize}
    \item Crescem exponencialmente sem predadores
    \item Perdem individuos por predação
    \item Taxa de predação $\propto$ encontros ($VP$)
\end{itemize}

\column{0.48\textwidth}
\textbf{Equação dos Predadores:}
\begin{equation*}
\derivada{P}{t} = \underbrace{\delta \beta V P}_{\text{crescimento}} - \underbrace{\gamma P}_{\text{morte}}
\end{equation*}

\begin{itemize}
    \item Crescem ao consumir presas
    \item Morrem naturalmente
    \item Dependem das presas para sobreviver
\end{itemize}
\end{columns}

\vspace{0.3cm}

\importante{
Este sistema exibe oscilações periódicas características!
}
\end{frame}

\begin{frame}[fragile, allowframebreaks]{Implementação: Lotka-Volterra}
\begin{lstlisting}[language=Python, caption=Sistema predador-presa, basicstyle=\ttfamily\footnotesize]
import numpy as np
import matplotlib.pyplot as plt
from scipy.integrate import odeint

def lotka_volterra(y, t, alpha, beta, delta, gamma):
    """Sistema Lotka-Volterra"""
    V, P = y
    dVdt = alpha*V - beta*V*P
    dPdt = delta*beta*V*P - gamma*P
    return [dVdt, dPdt]

# Parametros
alpha, beta, delta, gamma = 1.0, 0.1, 0.75, 1.5
y0, t = [10, 5], np.linspace(0, 50, 1000)
\end{lstlisting}
\end{frame}

\begin{frame}[fragile, shrink=5]{Implementação: Lotka-Volterra (2/2)}
\begin{lstlisting}[language=Python, caption=Resolução e Plotting]
# Resolver
sol = odeint(lotka_volterra, y0, t,
             args=(alpha, beta, delta, gamma))

# Plotar series temporais
plt.plot(t, sol[:, 0], 'b-', label='Presas')
plt.plot(t, sol[:, 1], 'r-', label='Predadores')
plt.xlabel('Tempo')
plt.ylabel('Populacao')
plt.legend()
\end{lstlisting}
\end{frame}

\begin{frame}[fragile, shrink=10]{Retrato de Fase}
\begin{columns}
\column{0.48\textwidth}
\begin{lstlisting}[language=Python, caption=Plano de fase, basicstyle=\ttfamily\tiny]
# Retrato de fase
plt.figure()
plt.plot(sol[:, 0], sol[:, 1],
         'g-', linewidth=2)
plt.plot(sol[0, 0], sol[0, 1],
         'go', markersize=10,
         label='Inicio')
plt.xlabel('Presas (V)')
plt.ylabel('Predadores (P)')
plt.title('Retrato de Fase')
plt.legend()
plt.grid(True)

# Campo vetorial
V = np.linspace(0, 20, 20)
P = np.linspace(0, 15, 20)
V_grid, P_grid = np.meshgrid(V, P)
dV, dP = lotka_volterra(
    [V_grid, P_grid], 0,
    alpha, beta, delta, gamma)
plt.quiver(V_grid, P_grid,
           dV, dP, alpha=0.5)
\end{lstlisting}

\column{0.48\textwidth}
\begin{center}
\begin{tikzpicture}[scale=0.8]
    \draw[->] (0,0) -- (4,0) node[right] {$V$};
    \draw[->] (0,0) -- (0,4) node[above] {$P$};

    % Trajetoria em espiral
    \draw[thick, verdeacido, domain=0:720, samples=200, smooth]
        plot ({2 + 1.5*cos(\x)/(1+\x/360)},
              {2 + 1.5*sin(\x)/(1+\x/360)});

    % Ponto de equilibrio
    \fill[red] (2,2) circle (2pt) node[right, font=\tiny] {Equilíbrio};

    \node[below] at (2,-0.2) {Presas};
    \node[left, rotate=90] at (-0.3,2) {Predadores};
\end{tikzpicture}
\end{center}

\textbf{Características:}
\begin{itemize}
    \item Trajetórias fechadas
    \item Oscilações periódicas
    \item Centro estável
\end{itemize}
\end{columns}
\end{frame}

% ============================================================================
% SEÇÃO 6: PONTOS DE EQUILÍBRIO
% ============================================================================
\section{Analise de Equilíbrio}

\begin{frame}[shrink=5]{Pontos de Equilíbrio (Steady States)}
\definicao{Ponto de Equilíbrio}{
Estado no qual o sistema não muda com o tempo:
\begin{equation*}
\derivada{\vect{x}}{t} = \vect{0} \quad \Rightarrow \quad \vect{f}(\vect{x}^*, t) = \vect{0}
\end{equation*}
}

\vspace{0.3cm}

\begin{exampleblock}{Exemplo: Crescimento Logístico}
\begin{equation*}
\derivada{N}{t} = rN\left(1 - \frac{N}{K}\right) = 0
\end{equation*}

Pontos de equilíbrio:
\begin{itemize}
    \item $N^*_1 = 0$ (extinção)
    \item $N^*_2 = K$ (capacidade de suporte)
\end{itemize}
\end{exampleblock}
\end{frame}

\begin{frame}{Estabilidade de Pontos de Equilíbrio}
\begin{block}{Classificação}
\begin{itemize}
    \item \textbf{Estável:} sistema retorna ao equilíbrio após perturbação
    \item \textbf{Instável:} sistema se afasta após perturbação
    \item \textbf{Marginalmente estável:} comportamento intermediário
\end{itemize}
\end{block}

\vspace{0.3cm}

\begin{center}
\begin{tikzpicture}[scale=0.7]
    % Estavel
    \begin{scope}
        \draw[->] (-1,0) -- (3,0) node[right] {$x$};
        \draw[->] (0,-1) -- (0,2) node[above] {$\dot{x}$};
        \draw[thick, azulescuro] (-0.5,1.5) parabola bend (1,0) (2.5,1.5);
        \fill[verdeacido] (1,0) circle (2pt) node[below, font=\tiny] {Estável};
        \draw[->, thick, red] (0.5,0.3) -- (1,0.1);
        \draw[->, thick, red] (1.5,0.3) -- (1,0.1);
    \end{scope}

    % Instavel
    \begin{scope}[xshift=5cm]
        \draw[->] (-1,0) -- (3,0) node[right] {$x$};
        \draw[->] (0,-1) -- (0,2) node[above] {$\dot{x}$};
        \draw[thick, azulescuro] (-0.5,-1.5) parabola[bend at end] (1,0);
        \draw[thick, azulescuro] (1,0) parabola (2.5,-1.5);
        \fill[red] (1,0) circle (2pt) node[below, font=\tiny] {Instável};
        \draw[->, thick, verdeacido] (0.5,-0.3) -- (0,-0.5);
        \draw[->, thick, verdeacido] (1.5,-0.3) -- (2,-0.5);
    \end{scope}
\end{tikzpicture}
\end{center}
\end{frame}

\begin{frame}{Analise de Estabilidade Linear}
\begin{block}{Linearização em Torno do Equilíbrio}
Para $\vect{x} = \vect{x}^* + \delta\vect{x}$ (pequena perturbação):
\begin{equation*}
\derivada{\delta\vect{x}}{t} = \mathbf{J}(\vect{x}^*) \delta\vect{x}
\end{equation*}

onde $\mathbf{J}$ é a matriz Jacobiana:
\begin{equation*}
J_{ij} = \parcial{f_i}{x_j}\Bigg|_{\vect{x}^*}
\end{equation*}
\end{block}

\begin{exampleblock}{Critério de Estabilidade}
O equilíbrio $\vect{x}^*$ é estável se todos os autovalores de $\mathbf{J}(\vect{x}^*)$ têm parte real negativa.
\end{exampleblock}
\end{frame}

\begin{frame}{Exemplo: Analise do Crescimento Logístico}
\begin{block}{Modelo}
\begin{equation*}
\derivada{N}{t} = f(N) = rN\left(1 - \frac{N}{K}\right)
\end{equation*}
\end{block}

\textbf{Jacobiano (1D):}
\begin{equation*}
J = \derivada{f}{N} = r\left(1 - \frac{2N}{K}\right)
\end{equation*}

\textbf{Nos pontos de equilíbrio:}
\begin{itemize}
    \item $N^* = 0$: $J(0) = r > 0$ $\Rightarrow$ \textcolor{red}{instável}
    \item $N^* = K$: $J(K) = -r < 0$ $\Rightarrow$ \textcolor{verdeacido}{estável}
\end{itemize}

\importante{
A populacao converge para a capacidade de suporte $K$
}
\end{frame}

<<<<<<< HEAD
\begin{frame}[fragile, shrink=10]{Analise de Estabilidade Computacional}
\begin{lstlisting}[language=Python, caption=Calculo do Jacobiano e Autovalores, basicstyle=\ttfamily\tiny]
=======
\begin{frame}[fragile, allowframebreaks]{Analise de Estabilidade Computacional}
\begin{lstlisting}[language=Python, caption=Analise de autovalores, basicstyle=\ttfamily\footnotesize]
>>>>>>> e775cf3 (docs: fix overflows and add computational examples to chapter 02)
import numpy as np
from scipy.linalg import eig

def jacobian_lotka_volterra(y, alpha, beta, delta, gamma):
    V, P = y
    return np.array([
        [alpha - beta*P, -beta*V],
        [delta*beta*P, delta*beta*V - gamma]
    ])

# Ponto de equilibrio e Jacobiano
V_eq, P_eq = gamma / (delta * beta), alpha / beta
J = jacobian_lotka_volterra([V_eq, P_eq], alpha, beta, delta, gamma)

# Autovalores
eigenvalues, _ = eig(J)
print(f"Autovalores: {eigenvalues}")
\end{lstlisting}
\end{frame}

\begin{frame}[fragile]{Analise de Estabilidade Computacional}
\begin{lstlisting}[language=Python, caption=Analise de autovalores]
# Ponto de equilibrio
V_eq = gamma / (delta * beta)
P_eq = alpha / beta

# Calcular Jacobiano no equilibrio
J = jacobian_lotka_volterra([V_eq, P_eq],
                             alpha, beta, delta, gamma)

# Autovalores
eigenvalues, eigenvectors = eig(J)
print(f"Autovalores: {eigenvalues}")
print(f"Partes reais: {eigenvalues.real}")
\end{lstlisting}
\end{frame}

% ============================================================================
% SEÇÃO 7: MÉTODOS NUMÉRICOS
% ============================================================================
\section{Métodos Numéricos}

\begin{frame}{Por que Métodos Numéricos?}
\begin{alertblock}{Limitações das Soluções Analíticas}
\begin{itemize}
    \item Maioria dos sistemas biológicos são não-lineares
    \item Poucos modelos têm solução analítica fechada
    \item Sistemas complexos requerem aproximação numérica
\end{itemize}
\end{alertblock}

\vspace{0.3cm}

\begin{block}{Métodos Numéricos Comuns}
\begin{itemize}
    \item \textbf{Euler:} simples, primeira ordem
    \item \textbf{Runge-Kutta (RK4):} preciso, quarta ordem
    \item \textbf{Métodos adaptativos:} passo variável (RK45, Dormand-Prince)
    \item \textbf{Métodos implícitos:} sistemas rígidos (BDF)
\end{itemize}
\end{block}
\end{frame}

\begin{frame}[allowframebreaks]{Método de Euler}
\begin{block}{Ideia Básica}
Aproximar derivada por diferença finita:
\begin{equation*}
\derivada{x}{t} \approx \frac{x(t+h) - x(t)}{h}
\end{equation*}
\end{block}

\destaque{
\begin{equation*}
x(t+h) = x(t) + h \cdot f(x(t), t)
\end{equation*}
}

\textbf{Algoritmo:}
\begin{enumerate}
    \item Partir de $x_0$ em $t_0$
    \item Calcular $x_1 = x_0 + h \cdot f(x_0, t_0)$
    \item Repetir: $x_{n+1} = x_n + h \cdot f(x_n, t_n)$
\end{enumerate}

\begin{exampleblock}{Erro}
Erro local: $O(h^2)$, Erro global: $O(h)$
\end{exampleblock}
\end{frame}

\begin{frame}[fragile, allowframebreaks]{Implementação do Método de Euler}
\begin{lstlisting}[language=Python, caption=Euler forward, basicstyle=\ttfamily\footnotesize]
def euler(f, x0, t):
    n = len(t)
    x = np.zeros((n, len(x0)))
    x[0] = x0
    for i in range(n-1):
        h = t[i+1] - t[i]
        x[i+1] = x[i] + h * np.array(f(x[i], t[i]))
    return x

# Exemplo
def exp_growth(x, t, r=0.5): return r * x
t = np.linspace(0, 10, 100)
x_euler = euler(lambda x,t: exp_growth(x,t), [1.0], t)
\end{lstlisting}
\end{frame}

\begin{frame}[allowframebreaks]{Método de Runge-Kutta de 4ª Ordem (RK4)}
\begin{block}{Ideia}
Usar múltiplas avaliações da derivada para melhor aproximação
\end{block}

\destaque{
\begin{align*}
k_1 &= f(x_n, t_n)\\
k_2 &= f(x_n + \frac{h}{2}k_1, t_n + \frac{h}{2})\\
k_3 &= f(x_n + \frac{h}{2}k_2, t_n + \frac{h}{2})\\
k_4 &= f(x_n + hk_3, t_n + h)\\
x_{n+1} &= x_n + \frac{h}{6}(k_1 + 2k_2 + 2k_3 + k_4)
\end{align*}
}

\begin{exampleblock}{Precisão}
Erro local: $O(h^5)$, Erro global: $O(h^4)$
\end{exampleblock}
\end{frame}

\begin{frame}[fragile, allowframebreaks]{Implementação do RK4}
\begin{lstlisting}[language=Python, caption=Runge-Kutta 4, basicstyle=\ttfamily\footnotesize]
def rk4(f, x0, t):
    n, x = len(t), np.zeros((len(t), len(x0)))
    x[0] = x0
    for i in range(n-1):
        h = t[i+1] - t[i]
        k1 = np.array(f(x[i], t[i]))
        k2 = np.array(f(x[i] + h*k1/2, t[i] + h/2))
        k3 = np.array(f(x[i] + h*k2/2, t[i] + h/2))
        k4 = np.array(f(x[i] + h*k3, t[i] + h))
        x[i+1] = x[i] + (h/6) * (k1 + 2*k2 + 2*k3 + k4)
    return x

# Usar
x_rk4 = rk4(lambda x,t: exp_growth(x,t), [1.0], t)
\end{lstlisting}
\end{frame}

% ============================================================================
% SEÇÃO 8: ANÁLISE DE SENSIBILIDADE
% ============================================================================
\section{Analise de Sensibilidade}

\begin{frame}[allowframebreaks]{Analise de Sensibilidade}
\definicao{Sensibilidade}{
Medida de como a saída de um modelo varia em resposta a mudanças nos parâmetros de entrada.
}

\vspace{0.3cm}

\begin{block}{Por que é importante?}
\begin{itemize}
    \item Identificar parâmetros críticos
    \item Priorizar experimentos para estimar parâmetros
    \item Avaliar robustez do modelo
    \item Simplificar modelos complexos
\end{itemize}
\end{block}

\begin{exampleblock}{Coeficiente de Sensibilidade}
\begin{equation*}
S_{x,p} = \parcial{x}{p} \cdot \frac{p}{x} = \frac{\parcial \ln x}{\parcial \ln p}
\end{equation*}

Sensibilidade relativa normalizada
\end{exampleblock}
\end{frame}

\begin{frame}[allowframebreaks]{Tipos de Analise de Sensibilidade}
\begin{block}{Local}
\begin{itemize}
    \item Variação em torno de um ponto nominal
    \item Derivadas parciais
    \item Rápida e analítica (quando possível)
\end{itemize}
\end{block}

\begin{block}{Global}
\begin{itemize}
    \item Exploração de todo o espaço de parâmetros
    \item Métodos de amostragem (Monte Carlo, Sobol)
    \item Computacionalmente intensiva
\end{itemize}
\end{block}

\begin{alertblock}{Trade-off}
Local: rápida mas limitada \\
Global: abrangente mas custosa
\end{alertblock}
\end{frame}

\begin{frame}[fragile, allowframebreaks]{Analise de Sensibilidade Local}
\begin{lstlisting}[language=Python, caption=Sensibilidade por diferenças finitas, basicstyle=\ttfamily\footnotesize]
def sensitivity_local(model, params, param_names,
                       perturbation=0.01):
    """
    Analise de sensibilidade local
    """
    baseline = model(params)
    sensitivities = {}
    for i, name in enumerate(param_names):
        params_pert = params.copy()
        params_pert[i] *= (1 + perturbation)
        output_pert = model(params_pert)
        sensitivities[name] = (output_pert - baseline) / baseline / perturbation
    return sensitivities

# Exemplo
sens = sensitivity_local(lambda p: p[1], [0.5, 100, 10], ['r', 'K', 'N0'])
\end{lstlisting}
\end{frame}

% ============================================================================
% SEÇÃO 9: AJUSTE DE MODELOS A DADOS
% ============================================================================
\section{Ajuste de Modelos a Dados}

\begin{frame}[allowframebreaks]{Estimação de Parâmetros}
\begin{block}{Problema}
Dados: $(t_1, y_1), (t_2, y_2), \ldots, (t_n, y_n)$\\
Modelo: $\hat{y}(t; \vect{p})$\\
Objetivo: Encontrar $\vect{p}^*$ que melhor ajusta os dados
\end{block}

\vspace{0.3cm}

\begin{block}{Função Objetivo}
Mínimos quadrados:
\begin{equation*}
\text{SSE}(\vect{p}) = \sum_{i=1}^{n} [y_i - \hat{y}(t_i; \vect{p})]^2
\end{equation*}

Objetivo: $\min_{\vect{p}} \text{SSE}(\vect{p})$
\end{block}

\begin{exampleblock}{Métodos de Otimização}
\begin{itemize}
    \item Levenberg-Marquardt
    \item Evolução Diferencial
    \item Algoritmos Genéticos
\end{itemize}
\end{exampleblock}
\end{frame}

\begin{frame}[fragile, allowframebreaks]{Ajuste de Curvas com scipy}
\begin{lstlisting}[language=Python, caption=Estimação de parâmetros, basicstyle=\ttfamily\footnotesize]
from scipy.optimize import curve_fit

# Dados experimentais
t_data = np.linspace(0, 10, 20)
N_data = 100/(1+9*np.exp(-0.5*t_data)) + np.random.normal(0, 2, 20)

# Modelo e Ajuste
def logistic_model(t, r, K): return K / (1 + (K/10 - 1)*np.exp(-r*t))
p_opt, _ = curve_fit(logistic_model, t_data, N_data, p0=[1.0, 50])

# Plotar
plt.plot(t_data, N_data, 'o', label='Dados')
plt.plot(t_data, logistic_model(t_data, *p_opt), '-', label='Ajuste')
\end{lstlisting}
\end{frame}

\begin{frame}[allowframebreaks]{Validação de Modelos}
\begin{block}{Métricas de Qualidade do Ajuste}
\begin{itemize}
    \item \textbf{SSE:} Soma dos erros quadrados
    \item \textbf{R$^2$:} Coeficiente de determinação (0 a 1)
    \item \textbf{RMSE:} Raiz do erro quadrático médio
    \item \textbf{AIC/BIC:} Critérios de informação (penalizam complexidade)
\end{itemize}
\end{block}

\begin{equation*}
R^2 = 1 - \frac{\sum (y_i - \hat{y}_i)^2}{\sum (y_i - \bar{y})^2}
\end{equation*}

\begin{alertblock}{Atenção!}
\begin{itemize}
    \item $R^2$ alto não garante bom modelo
    \item Validar com dados independentes
    \item Verificar plausibilidade biológica
\end{itemize}
\end{alertblock}
\end{frame}

% ============================================================================
% SEÇÃO 10: EXEMPLO INTEGRADO
% ============================================================================
\section{Estudo de Caso}

\begin{frame}[shrink=15]{Estudo de Caso: Dinâmica de Infecção Viral}
\begin{block}{Sistema}
Modelo SIR (Susceptible-Infected-Recovered):
\begin{align*}
\derivada{S}{t} &= -\beta \frac{SI}{N}\\
\derivada{I}{t} &= \beta \frac{SI}{N} - \gamma I\\
\derivada{R}{t} &= \gamma I
\end{align*}
\end{block}

\begin{block}{Número Básico de Reprodução ($R_0$)}
\begin{equation*}
R_0 = \frac{\beta}{\gamma}
\end{equation*}
\end{block}
\end{frame}

\begin{frame}{Estudo de Caso: Dinâmica de Infecção Viral (cont.)}
\begin{block}{Parâmetros do Modelo SIR}
onde:
\begin{itemize}
    \item $S$ = susceptíveis
    \item $I$ = infectados
    \item $R$ = recuperados
    \item $\beta$ = taxa de transmissão
    \item $\gamma$ = taxa de recuperação
    \item $N = S + I + R$ = populacao total
\end{itemize}
\end{block}
\end{frame}

\begin{frame}[allowframebreaks]{Numero Básico de Reprodução}
\definicao{$R_0$ (R-naught)}{
Numero médio de infecções secundárias causadas por um indivíduo infectado em uma populacao totalmente susceptível.
}

\vspace{0.3cm}

\destaque{
\begin{equation*}
R_0 = \frac{\beta}{\gamma}
\end{equation*}
}

\begin{block}{Interpretação}
\begin{itemize}
    \item $R_0 < 1$: epidemia se extingue
    \item $R_0 > 1$: epidemia se propaga
    \item $R_0 = 1$: limiar epidêmico
\end{itemize}
\end{block}

\exemplo{COVID-19}{
$R_0$ estimado entre 2-6, dependendo da variante e medidas de controle
}
\end{frame}

\begin{frame}[fragile, allowframebreaks]{Implementação do Modelo SIR}
\begin{lstlisting}[language=Python, caption=Modelo SIR completo, basicstyle=\ttfamily\footnotesize]
def sir_model(y, t, beta, gamma):
    S, I, R = y
    N = S + I + R
    return [-beta * S * I / N, beta * S * I / N - gamma * I, gamma * I]

# Parametros e Condições
beta, gamma = 0.5, 0.1
R0 = beta / gamma
N, I0 = 1000, 10
y0 = [N - I0, I0, 0]
\end{lstlisting}
\end{frame}

\begin{frame}[fragile, allowframebreaks]{Analise e Visualização}
\begin{lstlisting}[language=Python, caption=Analise do modelo SIR, basicstyle=\ttfamily\footnotesize]
S, I, R = sol.T
plt.figure(figsize=(10, 6))
plt.plot(t, S, 'b-', label='Suscept.'), plt.plot(t, I, 'r-', label='Infect.')
plt.plot(t, R, 'g-', label='Recuper.')
plt.xlabel('Tempo'), plt.ylabel('Individuos'), plt.legend(), plt.grid(True)

# Estatisticas
peak_idx = np.argmax(I)
print(f"Pico: {I[peak_idx]:.0f} em {t[peak_idx]:.1f} dias")
print(f"Ataque final: {R[-1]/N*100:.1f}%")
\end{lstlisting}
\end{frame}

% ============================================================================
% SEÇÃO 11: BOAS PRÁTICAS
% ============================================================================
\section{Boas Práticas em Modelagem}

\begin{frame}[shrink=15]{Princípios de Modelagem}
\begin{block}{1. Simplicidade vs Realismo}
\begin{itemize}
    \item Comece simples, adicione complexidade gradualmente
    \item "Make things as simple as possible, but not simpler" --- Einstein
    \item Balanceie entre descrição acurada e tratabilidade
\end{itemize}
\end{block}

\begin{block}{2. Validação}
\begin{itemize}
    \item Compare com dados experimentais
    \item Teste previsões em novos experimentos
    \item Valide em múltiplas escalas/condições
\end{itemize}
\end{block}

\begin{block}{3. Documentação}
\begin{itemize}
    \item Documente hipóteses claramente
    \item Registre origem de parâmetros
    \item Mantenha código reproduzível
\end{itemize}
\end{block}
\end{frame}

\begin{frame}[shrink]{Armadilhas Comuns}
\begin{alertblock}{Erros Frequentes}
\begin{enumerate}
    \item \textbf{Overfitting:} modelo complexo demais para os dados
    \item \textbf{Parâmetros não-identificáveis:} múltiplas soluções
    \item \textbf{Unidades inconsistentes:} sempre verifique dimensões
    \item \textbf{Condições iniciais irrealistas:} valores negativos, etc.
    \item \textbf{Extrapolação excessiva:} além do regime validado
    \item \textbf{Ignorar incerteza:} sempre quantifique incerteza
\end{enumerate}
\end{alertblock}

\importante{
"Remember that all models are wrong; the practical question is how wrong do they have to be to not be useful" --- George Box
}
\end{frame}

\begin{frame}{Ferramentas Computacionais}
\begin{columns}
\column{0.48\textwidth}
\begin{block}{Python}
\begin{itemize}
    \item NumPy, SciPy
    \item Matplotlib, Seaborn
    \item PyDSTool
    \item Tellurium
    \item COBRApy (metabolismo)
\end{itemize}
\end{block}

\begin{block}{R}
\begin{itemize}
    \item deSolve
    \item FME (fitting)
    \item ggplot2
\end{itemize}
\end{block}

\column{0.48\textwidth}
\begin{block}{MATLAB}
\begin{itemize}
    \item ode45, ode15s
    \item SimBiology
    \item Simbólico
\end{itemize}
\end{block}

\begin{block}{Especializadas}
\begin{itemize}
    \item COPASI
    \item CellDesigner
    \item Virtual Cell
    \item BioNetGen
\end{itemize}
\end{block}
\end{columns}
\end{frame}

% ============================================================================
% SEÇÃO 12: EXERCÍCIOS
% ============================================================================
\section{Exercícios}

\begin{frame}[allowframebreaks]{Exercícios - Parte 1}
\begin{block}{Questão 1: Modelo de Crescimento}
Considere um modelo de crescimento com colheita constante:
\begin{equation*}
\derivada{N}{t} = rN\left(1 - \frac{N}{K}\right) - H
\end{equation*}
\begin{enumerate}
    \item Encontre os pontos de equilíbrio
    \item Analise a estabilidade para $H < rK/4$ e $H > rK/4$
\end{enumerate}
\end{block}
\end{frame}

\begin{frame}{Exercícios - Parte 2}
\begin{block}{Questão 2: Método Numérico}
Implemente o método de Euler para resolver:
\begin{equation*}
\derivada{y}{t} = -y + \sin(t), \quad y(0) = 1
\end{equation*}
Compare com a solução analítica para diferentes passos.
\end{block}
\end{frame}

\begin{frame}{Exercícios - Parte 2}
\begin{block}{Questão 3: Sistema Acoplado}
Considere o sistema de competição:
\begin{align*}
\derivada{N_1}{t} &= r_1 N_1 \left(1 - \frac{N_1 + \alpha N_2}{K_1}\right)\\
\derivada{N_2}{t} &= r_2 N_2 \left(1 - \frac{N_2 + \beta N_1}{K_2}\right)
\end{align*}

\begin{enumerate}
    \item Implemente o modelo em Python
    \item Encontre os pontos de equilíbrio
    \item Simule para diferentes valores de $\alpha$ e $\beta$
    \item O que acontece quando $\alpha\beta > 1$ vs $\alpha\beta < 1$?
\end{enumerate}
\end{block}
\end{frame}

\begin{frame}[shrink=20]{Exercícios - Parte 3}
\begin{alertblock}{Projeto: Modelo SIR com Vacinação}
Estenda o modelo SIR para incluir vacinação:
\begin{align*}
\derivada{S}{t} &= -\beta \frac{SI}{N} - \nu S, \quad \derivada{I}{t} = \beta \frac{SI}{N} - \gamma I, \quad \derivada{R}{t} = \gamma I + \nu S
\end{align*}
\begin{enumerate}
    \item Implemente o modelo e calcule o $R_0$ efetivo
    \item Determine $\nu$ necessária para controle ($R_0^{eff} < 1$)
    \item Simule diferentes estratégias de vacinação
\end{enumerate}
\end{alertblock}
\end{frame}

% ============================================================================
% SEÇÃO 13: REFERÊNCIAS
% ============================================================================
\section{Referências}

\begin{frame}{Referências Principais}
\begin{thebibliography}{99}
\bibitem{strogatz} Strogatz S.H. (2018). \textit{Nonlinear Dynamics and Chaos}, 2nd ed. CRC Press.

\bibitem{edelstein} Edelstein-Keshet L. (2005). \textit{Mathematical Models in Biology}. SIAM.

\bibitem{britton} Britton N.F. (2003). \textit{Essential Mathematical Biology}. Springer.

\bibitem{murray} Murray J.D. (2002). \textit{Mathematical Biology I: An Introduction}, 3rd ed. Springer.

\bibitem{keener} Keener J., Sneyd J. (2009). \textit{Mathematical Physiology}, 2nd ed. Springer.
\end{thebibliography}
\end{frame}

\begin{frame}[allowframebreaks]{Leitura Complementar}
\begin{block}{Livros Avançados}
\begin{itemize}
    \item Ellner S.P., Guckenheimer J. (2011). \textit{Dynamic Models in Biology}
    \item Alon U. (2019). \textit{An Introduction to Systems Biology}, 2nd ed.
    \item Klipp E. et al. (2016). \textit{Systems Biology}, 2nd ed.
\end{itemize}
\end{block}

\begin{block}{Recursos Online}
\begin{itemize}
    \item \textbf{Wolfram MathWorld:} \url{https://mathworld.wolfram.com}
    \item \textbf{MIT OCW:} Mathematical Biology course
    \item \textbf{Scipy Lecture Notes:} \url{https://scipy-lectures.org}
\end{itemize}
\end{block}

\begin{block}{Tutoriais Python}
\begin{itemize}
    \item \url{https://scipy-cookbook.readthedocs.io}
    \item \url{https://python-numerics.github.io}
\end{itemize}
\end{block}
\end{frame}

% ============================================================================
% SLIDE FINAL
% ============================================================================
\begin{frame}
\begin{center}
{\Huge Obrigado!}

\vspace{1cm}

{\Large Capítulo 2: Introdução à Modelagem Matemática}

\vspace{1cm}

\textbf{Próximo Capítulo:}\\
Static Network Models

\vspace{1cm}

\textit{Dúvidas? Perguntas?}
\end{center}
\end{frame}

\end{document}
